\documentclass[12pt]{article}
\usepackage{../../ficha}

\begin{document}

\ficha{Limites de Sucessões e o Número $e$}

\instrucoes{Resolve os exercícios, apresentando os cálculos e levantando as indeterminações quando necessário.}

\secao{Noção de Limite}{Resolve.}
\begin{exercicios}
  \item Calcula os termos de ordem 10, 100 e 1\,000 de $u_n = \dfrac{1}{n}$. Para que valor parecem tender?
  \item Determina a partir de que ordem os termos de $u_n = \dfrac{1}{n}$ são menores do que $0{,}001$.
  \item Mostra que $u_n = \dfrac{2n + 1}{n}$ tende para 2, determinando a partir de que ordem $|u_n - 2| < 0{,}01$.
  \item Dá um exemplo de uma sucessão que tende para $+\infty$, uma que tende para $-\infty$ e uma que não tem limite.
  \item Uma sucessão limitada é sempre convergente? Dá um exemplo.
  \item Uma sucessão monótona e limitada é convergente? Que teorema garante isso?
\end{exercicios}

\secao{Limites de Sucessões Polinomiais e Racionais}{Calcula o limite.}
\begin{exercicios}
  \item $\lim (3n^2 - 5n + 1)$
  \item $\lim (-n^3 + 100n^2)$
  \item $\lim \dfrac{4n + 1}{2n - 3}$
  \item $\lim \dfrac{n^2 + 1}{3n^2 - n}$
  \item $\lim \dfrac{5n + 2}{n^2 + 1}$
  \item $\lim \dfrac{n^3 - 2}{4n^2 + n}$
  \item $\lim \dfrac{(2n + 1)^2}{n^2 + 5}$
  \item $\lim \dfrac{1 - 2n^2}{n + 3}$
  \item $\lim \dfrac{3 - n}{n + 1}$
  \item $\lim \left(\dfrac{1}{n} + 5\right)\left(2 - \dfrac{3}{n^2}\right)$
\end{exercicios}

\secao{Indeterminações com Radicais}{Calcula o limite.}
\begin{exercicios}
  \item $\lim \left(\sqrt{n + 1} - \sqrt{n}\right)$
  \item $\lim \left(\sqrt{n^2 + 3n} - n\right)$
  \item $\lim \left(\sqrt{4n^2 + 1} - 2n\right)$
  \item $\lim \dfrac{\sqrt{n^2 + 1}}{n}$
  \item $\lim \dfrac{\sqrt{9n^2 + n}}{2n + 1}$
  \item $\lim \dfrac{n}{\sqrt{n + 1}}$
\end{exercicios}

\secao{Sucessões Exponenciais e Progressões}{Calcula o limite.}
\begin{exercicios}
  \item $\lim 2^n$
  \item $\lim \left(\dfrac{1}{3}\right)^n$
  \item $\lim (-2)^n$
  \item $\lim \dfrac{3^n + 1}{3^n}$
  \item $\lim \dfrac{2^n + 5^n}{5^n - 1}$
  \item $\lim \dfrac{4^n}{2^{2n+1}}$
  \item $\lim \left(1 + \dfrac{1}{2} + \dfrac{1}{4} + \cdots + \dfrac{1}{2^{n-1}}\right)$
  \item Calcula a soma de todos os termos da progressão geométrica $9, 3, 1, \dfrac{1}{3}, \ldots$
\end{exercicios}

\secao{O Número $e$}{Usa $\lim \left(1 + \dfrac{1}{n}\right)^n = e$ e $\lim \left(1 + \dfrac{x}{n}\right)^n = e^x$.}
\begin{exercicios}
  \item $\lim \left(1 + \dfrac{1}{n}\right)^{2n}$
  \item $\lim \left(1 + \dfrac{3}{n}\right)^n$
  \item $\lim \left(1 - \dfrac{2}{n}\right)^n$
  \item $\lim \left(1 + \dfrac{1}{n}\right)^{n+5}$
  \item $\lim \left(\dfrac{n + 2}{n}\right)^n$
  \item $\lim \left(\dfrac{n + 1}{n - 1}\right)^n$
  \item $\lim \left(\dfrac{n}{n + 3}\right)^{2n}$
  \item Calcula, com calculadora, $\left(1 + \dfrac{1}{n}\right)^n$ para $n = 10$, $100$ e $10\,000$. Compara com $e \approx 2{,}71828$.
\end{exercicios}

\secao{Desafios}{Resolve os problemas.}
\begin{exercicios}
  \item Um capital de 1\,000 euros é aplicado à taxa anual de 6\,\%, com capitalização $n$ vezes por ano. Escreve o capital ao fim de um ano e calcula o seu limite quando $n \to +\infty$ (capitalização contínua).
  \item Usa o teorema das sucessões enquadradas para calcular $\lim \dfrac{\sin n}{n}$.
  \item Usa o teorema das sucessões enquadradas para calcular $\lim \dfrac{(-1)^n + n}{n}$.
  \item Mostra que a sucessão $u_1 = 1$, $u_{n+1} = \sqrt{2 + u_n}$ é crescente e majorada por 2. Determina o seu limite.
  \item Calcula $\lim \dfrac{1 + 2 + 3 + \cdots + n}{n^2}$.
\end{exercicios}

\end{document}
